\documentclass[11pt,a4paper,fontset=none]{ctexart}
\usepackage{fontspec}
\usepackage{xeCJK}
\usepackage{amsmath,amssymb,mathtools}
\usepackage{geometry}
\usepackage[dvipsnames]{xcolor}
\usepackage{tikz}
\usetikzlibrary{
  arrows.meta,
  angles,
  backgrounds,
  calc,
  decorations.markings,
  decorations.pathmorphing,
  decorations.pathreplacing,
  fit,
  intersections,
  matrix,
  patterns,
  patterns.meta,
  positioning,
  quotes,
  shapes.geometric,
  shapes.misc,
  through,
  3d
}
\usepackage{pgfplots}
\pgfplotsset{compat=1.18}
\usepackage[most]{tcolorbox}
\usepackage{enumitem}
\usepackage{multicol}
\usepackage{tabularx,array,booktabs,multirow}
\usepackage{fancyhdr}
\usepackage{lastpage}
\usepackage{titlesec}
\usepackage{needspace}
\usepackage{eso-pic}
\usepackage{graphicx}
\usepackage{hyperref}
\usepackage{xparse}

\geometry{
  a4paper,
  left=18mm,
  right=18mm,
  top=22mm,
  bottom=20mm,
  headheight=22pt,
  headsep=8mm,
  footskip=11mm
}

\setmainfont{texgyrepagella-regular.otf}[
  BoldFont=texgyrepagella-bold.otf,
  ItalicFont=texgyrepagella-italic.otf,
  BoldItalicFont=texgyrepagella-bolditalic.otf
]
\setsansfont{texgyreheros-regular.otf}[
  BoldFont=texgyreheros-bold.otf,
  ItalicFont=texgyreheros-italic.otf,
  BoldItalicFont=texgyreheros-bolditalic.otf
]
\setmonofont{texgyrecursor-regular.otf}[
  BoldFont=texgyrecursor-bold.otf
]
\setCJKmainfont{Kaiti SC}[
  BoldFont=Kaiti SC
]
\setCJKsansfont{FandolHei-Regular.otf}[
  BoldFont=FandolHei-Bold.otf
]
\setCJKmonofont{FandolKai-Regular.otf}

\definecolor{Paper}{HTML}{FBFCFA}
\definecolor{Ink}{HTML}{17324D}
\definecolor{Muted}{HTML}{647784}
\definecolor{Teal}{HTML}{167D89}
\definecolor{TealDark}{HTML}{0F5F69}
\definecolor{Mint}{HTML}{E8F5F2}
\definecolor{Sky}{HTML}{EAF3F8}
\definecolor{Coral}{HTML}{E66A5A}
\definecolor{Gold}{HTML}{D7A94A}
\definecolor{Rule}{HTML}{C8D8DA}
\definecolor{SolutionBg}{HTML}{FFF9F1}

\pagecolor{Paper}
\color{Ink}
\linespread{1.14}
\setlength{\parindent}{2em}
\setlength{\parskip}{0.25em}
\setlength{\columnsep}{1.3cm}
\setlength{\abovedisplayskip}{0.7em}
\setlength{\belowdisplayskip}{0.7em}
\setlist{nosep}

\hypersetup{
  colorlinks=true,
  linkcolor=TealDark,
  urlcolor=TealDark,
  pdfborder={0 0 0}
}

\newcommand{\papertitle}{Mathstalents 小学版练习题}
\newcommand{\SetPaperTitle}[1]{\renewcommand{\papertitle}{#1}}

\pagestyle{fancy}
\fancyhf{}
\fancyhead[L]{\sffamily\small\color{TealDark}\papertitle}
\fancyhead[R]{}
\fancyfoot[C]{%
  {\sffamily\scriptsize\color{Muted}第~\arabic{page}~页~/~共~\pageref*{LastPage}~页}%
}
\renewcommand{\headrulewidth}{0.4pt}
\renewcommand{\footrulewidth}{0pt}
\renewcommand{\headrule}{\hbox to\headwidth{\color{Rule}\leaders\hrule height \headrulewidth\hfill}}

\AddToShipoutPictureBG{%
  \begin{tikzpicture}[remember picture,overlay]
    \draw[Rule,line width=0.35pt,rounded corners=2pt]
      ([xshift=6mm,yshift=-6mm]current page.north west)
      rectangle
      ([xshift=-6mm,yshift=6mm]current page.south east);
    \draw[Teal,line width=1.1pt]
      ([xshift=6mm,yshift=-14mm]current page.north west) --
      ([xshift=6mm,yshift=-6mm]current page.north west) --
      ([xshift=14mm,yshift=-6mm]current page.north west);
    \draw[Coral,line width=1.1pt]
      ([xshift=-6mm,yshift=14mm]current page.south east) --
      ([xshift=-6mm,yshift=6mm]current page.south east) --
      ([xshift=-14mm,yshift=6mm]current page.south east);
  \end{tikzpicture}%
}

\newcommand{\metachip}[2][Teal]{%
  \tikz[baseline=(chip.base)]\node[
    fill=#1!10,
    draw=#1!45,
    text=#1!80!black,
    rounded corners=2.5pt,
    inner xsep=6pt,
    inner ysep=2.4pt,
    font=\sffamily\scriptsize
  ] (chip) {\textbf{#2}};%
}

\NewDocumentCommand{\makecover}{m m m m m m}{%
  \thispagestyle{empty}
  \begin{tikzpicture}[remember picture,overlay]
    \draw[Teal!14,line width=0.8pt,step=10mm]
      ([xshift=8mm,yshift=8mm]current page.south west) grid
      ([xshift=-8mm,yshift=-8mm]current page.north east);
    \node[
      anchor=north west,
      text=Ink,
      align=left,
      text width=155mm
    ] at ([xshift=24mm,yshift=-35mm]current page.north west) {
      {\sffamily\bfseries\fontsize{14}{18}\selectfont
        MATHSTALENTS \(\cdot\) ELEMENTARY}\\[9mm]
      {\sffamily\bfseries\fontsize{34}{42}\selectfont
        小学版 \(\,\boldsymbol{\cdot}\,\) #1}\\[3mm]
      {\sffamily\fontsize{17}{23}\selectfont
        练习题与参考答案}\\[11mm]
      {\color{TealDark}\rule{112mm}{1.1pt}}\\[7mm]
      {\sffamily\fontsize{13}{20}\selectfont #2}
    };
    \node[
      anchor=north west,
      fill=white,
      draw=TealDark,
      line width=0.6pt,
      text opacity=1,
      rounded corners=5pt,
      inner sep=9mm,
      text width=137mm
    ] at ([xshift=24mm,yshift=-126mm]current page.north west) {
      \begin{minipage}{137mm}
        \sffamily
        \begin{tabularx}{\linewidth}{>{\centering\arraybackslash}X
          >{\centering\arraybackslash}X
          >{\centering\arraybackslash}X}
          {\color{TealDark}\bfseries\Large #3} &
          {\color{TealDark}\bfseries\Large #4} &
          {\color{TealDark}\bfseries\Large 精编}\\[-1mm]
          {\color{Muted}\small 总分} &
          {\color{Muted}\small 题目数} &
          {\color{Muted}\small 题解版本}
        \end{tabularx}
        \vspace{6mm}
        {\color{Rule}\hrule}
        \vspace{5mm}
        {\color{Ink}\bfseries 题型结构}\par
        \vspace{2mm}
        {\color{Muted}\small #5}
      \end{minipage}
    };
  \end{tikzpicture}
  \null
  \clearpage
}

\newcommand{\paperinfo}[3]{%
  \begin{tcolorbox}[
    enhanced,
    colback=Sky,
    colframe=Teal!40,
    boxrule=0.6pt,
    arc=3pt,
    left=5mm,right=5mm,top=4mm,bottom=4mm
  ]
    \sffamily
    \begin{tabularx}{\linewidth}{@{}X X X@{}}
      \textbf{姓名：}\rule{28mm}{0.4pt} &
      \textbf{用时：}\rule{22mm}{0.4pt} &
      \textbf{得分：}\rule{22mm}{0.4pt}\\[3mm]
      \multicolumn{3}{@{}l@{}}{\small\color{Muted}
        #1 \qquad #2 \qquad #3}
    \end{tabularx}
  \end{tcolorbox}
  \vspace{2mm}
}

\newcommand{\mtQuestionNeedspace}{8.2cm}
\newcommand{\mtBannerNeedspace}{6.8cm}
\newcommand{\mtQuestionBreakLines}{8}

\newcommand{\sectionbanner}[2]{%
  \Needspace{\mtBannerNeedspace}
  \begin{tcolorbox}[
    enhanced,
    colback=white,
    colframe=TealDark,
    boxrule=0.65pt,
    arc=2.5pt,
    left=4mm,right=4mm,top=2.6mm,bottom=2.6mm,
    before skip=3.2mm,
    after skip=2.2mm
  ]
    {\sffamily\bfseries\large\color{TealDark}#1\quad
     \sffamily\small #2}
  \end{tcolorbox}
}

\NewDocumentEnvironment{question}{m m}{%
  \Needspace{\mtQuestionNeedspace}
  \begin{tcolorbox}[
    enhanced,
    breakable,
    lines before break=\mtQuestionBreakLines,
    colback=white,
    colframe=Rule,
    boxrule=0.55pt,
    arc=3pt,
    left=5mm,right=5mm,top=5mm,bottom=4mm,
    before skip=4mm,
    after skip=4mm,
    title={\sffamily\bfseries 第 #1 题（ #2 分 ）},
    coltitle=TealDark,
    colbacktitle=white,
    fonttitle=\sffamily\bfseries,
    attach boxed title to top left={xshift=4mm,yshift=-2.8mm},
    boxed title style={
      colback=white,
      colframe=TealDark,
      arc=2.5pt,
      boxrule=0.65pt,
      left=3mm,right=3mm,top=1.5mm,bottom=1.5mm
    }
  ]
  \hypertarget{question:#1}{}%
  \vspace{1mm}
}{%
  \end{tcolorbox}
}

\ExplSyntaxOn
\seq_new:N \g_mt_choices_seq
\dim_new:N \l_mt_choice_max_dim
\dim_new:N \l_mt_choice_sum_dim
\dim_new:N \l_mt_choice_tmp_dim
\dim_new:N \l_mt_choice_label_dim
\box_new:N \l_mt_choice_box

\NewDocumentCommand{\choice}{m}
{
  \seq_gput_right:Nn \g_mt_choices_seq {#1}
}

\cs_new_protected:Npn \mt_set_choice_label_dim:
{
  \hbox_set:Nn \l_mt_choice_box
  {
    \textcolor{TealDark}{\bfseries E.}
    \hspace{0.45em}
  }
  \dim_set:Nn \l_mt_choice_label_dim { \box_wd:N \l_mt_choice_box }
}

\cs_new_protected:Npn \mt_choice_label:n #1
{
  \textcolor{TealDark}{\bfseries \int_to_Alph:n {#1}.}
}

\cs_new_protected:Npn \mt_typeset_choice:nn #1#2
{
  \mt_choice_label:n {#1}\hspace{0.45em}#2
}

\cs_new_protected:Npn \mt_typeset_choice_par:nn #1#2
{
  \par\noindent
  \hangindent = \l_mt_choice_label_dim
  \hangafter = 1
  \makebox[\l_mt_choice_label_dim][l]{ \mt_choice_label:n {#1} }
  #2
  \par
}

\cs_new_protected:Npn \mt_measure_choices:
{
  \dim_zero:N \l_mt_choice_max_dim
  \dim_zero:N \l_mt_choice_sum_dim
  \int_step_inline:nn { \seq_count:N \g_mt_choices_seq }
  {
    \hbox_set:Nn \l_mt_choice_box
    {
      \mt_typeset_choice:nn {##1} { \seq_item:Nn \g_mt_choices_seq {##1} }
    }
    \dim_set:Nn \l_mt_choice_tmp_dim { \box_wd:N \l_mt_choice_box }
    \dim_set:Nn \l_mt_choice_sum_dim { \l_mt_choice_sum_dim + \l_mt_choice_tmp_dim + 1.6em }
    \dim_set:Nn \l_mt_choice_max_dim
    {
      \dim_max:nn { \l_mt_choice_max_dim } { \l_mt_choice_tmp_dim }
    }
  }
}

\cs_new_protected:Npn \mt_print_choices:
{
  \mt_set_choice_label_dim:
  \mt_measure_choices:
  \par\nopagebreak\vspace{0.5em}
  \dim_compare:nNnTF { \l_mt_choice_sum_dim } < { \linewidth }
  {
    \noindent
    \int_step_inline:nn { \seq_count:N \g_mt_choices_seq }
    {
      \mbox
      {
        \mt_typeset_choice:nn {##1} { \seq_item:Nn \g_mt_choices_seq {##1} }
      }
      \hfill
    }
    \par
  }
  {
    \dim_compare:nNnTF { \l_mt_choice_max_dim + 0.6em } < { 0.5\linewidth }
    {
      \noindent
      \int_step_inline:nn { \seq_count:N \g_mt_choices_seq }
      {
        \begin{minipage}[t]{0.48\linewidth}
          \raggedright
          \mt_typeset_choice_par:nn {##1} { \seq_item:Nn \g_mt_choices_seq {##1} }
        \end{minipage}
        \int_if_odd:nTF {##1}
        { \hfill }
        { \par\vspace{0.38em} }
      }
      \int_if_odd:nT { \seq_count:N \g_mt_choices_seq } { \par }
    }
    {
      \int_step_inline:nn { \seq_count:N \g_mt_choices_seq }
      {
        \mt_typeset_choice_par:nn {##1} { \seq_item:Nn \g_mt_choices_seq {##1} }
        \vspace{0.28em}
      }
    }
  }
  \vspace{0.25em}
  \seq_gclear:N \g_mt_choices_seq
}

\NewDocumentEnvironment{choices}{}
{
  \seq_gclear:N \g_mt_choices_seq
}
{
  \mt_print_choices:
}

\NewDocumentEnvironment{multichoices}{}
{
  \seq_gclear:N \g_mt_choices_seq
}
{
  \mt_print_choices:
}
\ExplSyntaxOff

\NewDocumentCommand{\subquestion}{m m}{%
  \par\medskip
  \noindent
  \metachip[Coral]{#1 分}\hspace{0.7em}#2\par
}

\newcommand{\answerlineskip}{9mm}

\newcommand{\answerspace}{%
  \par\vspace{12mm}%
}

\newcommand{\answerline}{%
  \ifmmode
    \underline{\hspace{2.8cm}}%
  \else
    \ifvmode
      \par\vspace{\answerlineskip}%
      \noindent{\color{Ink}\rule{0.46\linewidth}{0.5pt}}%
      \par\vspace{2.5mm}%
    \else
      \underline{\hspace{2.8cm}}%
    \fi
  \fi
}

\NewDocumentCommand{\answerbox}{m}{%
  \par\vspace{3mm}
  \begin{tcolorbox}[
    enhanced,
    colback=white,
    colframe=Teal!38!Rule,
    boxrule=0.7pt,
    arc=2.5pt,
    left=3.5mm,
    right=3.5mm,
    top=3mm,
    bottom=3mm,
    height=#1,
    before skip=0pt,
    after skip=2mm
  ]
    \mbox{}
  \end{tcolorbox}
}

\NewDocumentEnvironment{solution}{m m m m}{%
  \Needspace{5.4cm}
  \begin{tcolorbox}[
    enhanced,
    breakable,
    lines before break=6,
    colback=white,
    colframe=Gold!65!Ink,
    boxrule=0.6pt,
    arc=3pt,
    left=5mm,right=5mm,top=5mm,bottom=4mm,
    before skip=4mm,
    after skip=4mm,
    title={\sffamily\bfseries 第 #1 题 \(\cdot\) 参考解答},
    coltitle=Ink,
    colbacktitle=white,
    fonttitle=\sffamily\bfseries,
    attach boxed title to top left={xshift=4mm,yshift=-2.8mm},
    boxed title style={
      colback=white,
      colframe=Gold!65!Ink,
      arc=2.5pt,
      boxrule=0.65pt,
      left=3mm,right=3mm,top=1.5mm,bottom=1.5mm
    }
  ]
  \hypertarget{solution:#1}{}%
  \vspace{1mm}
  \noindent
  \metachip[Coral]{答案：#2}\hspace{0.45em}
  \metachip[Gold]{难度：#3}\hspace{0.45em}
  \metachip[Teal]{#4}
  \par\medskip
}{%
  \end{tcolorbox}
}

\NewDocumentCommand{\analysis}{+m}{%
  \par\smallskip
  {\sffamily\bfseries\color{TealDark}思路分析}\par
  \vspace{1mm}
  #1
  \par\smallskip
}

\NewDocumentCommand{\methodsummary}{+m}{%
  \par\medskip
  \noindent
  {\sffamily\bfseries\color{TealDark}方法总结}\quad
  {\sffamily\color{Ink}#1}\par
}

\newcommand{\solutiondivider}{%
  \clearpage
}


\SetPaperTitle{Mathstalents 小学版 · 知识点纲要}
\hypersetup{
  pdftitle={Mathstalents 小学版知识点纲要},
  pdfauthor={根据第三、四、五套卷整理}
}
\setcounter{tocdepth}{2}
\setlist[itemize]{leftmargin=1.55em,itemsep=0.22em,topsep=0.4em,parsep=0pt}

\newcommand{\outlinecover}{%
  \thispagestyle{empty}%
  \begin{tikzpicture}[remember picture,overlay]
    \draw[Teal!14,line width=0.8pt,step=10mm]
      ([xshift=8mm,yshift=8mm]current page.south west) grid
      ([xshift=-8mm,yshift=-8mm]current page.north east);
    \node[
      anchor=north west,
      text=Ink,
      align=left,
      text width=155mm
    ] at ([xshift=24mm,yshift=-32mm]current page.north west) {
      {\sffamily\bfseries\fontsize{14}{18}\selectfont
        MATHSTALENTS \(\cdot\) ELEMENTARY}\\[8mm]
      {\sffamily\bfseries\fontsize{32}{40}\selectfont
        小学版 \(\,\boldsymbol{\cdot}\,\) 知识点纲要}\\[4mm]
      {\sffamily\fontsize{16}{22}\selectfont
        第三、四、五套卷的基础工具}\\[10mm]
      {\color{TealDark}\rule{112mm}{1.1pt}}\\[6mm]
      {\sffamily\fontsize{13}{20}\selectfont 2026 年 9 月}
    };
    \node[
      anchor=north west,
      fill=white,
      draw=TealDark,
      line width=0.6pt,
      rounded corners=5pt,
      inner sep=8mm,
      text width=145mm
    ] at ([xshift=24mm,yshift=-118mm]current page.north west) {
      \begin{minipage}{145mm}
        \sffamily
        \begin{tabularx}{\linewidth}{>{\centering\arraybackslash}X
          >{\centering\arraybackslash}X
          >{\centering\arraybackslash}X}
          {\color{TealDark}\bfseries\Large 7} &
          {\color{TealDark}\bfseries\Large 3} &
          {\color{TealDark}\bfseries\Large 纲要}\\[-1mm]
          {\color{Muted}\small 知识类别} &
          {\color{Muted}\small 依据套卷} &
          {\color{Muted}\small 整理方式}
        \end{tabularx}
        \vspace{5.5mm}
        {\color{Rule}\hrule}
        \vspace{4.5mm}
        {\color{Ink}\bfseries 类别}\par
        \vspace{1.6mm}
        {\color{Muted}\small
          计数 \quad 概率 \quad 递推与周期 \quad 整除与因数\par
          \vspace{1.1mm}
          等式与最值 \quad 图形与运动 \quad 操作与不变量\par}
        \vspace{3.4mm}
        {\color{Ink}\small
          前七章写套卷里的工具：定义、一条事实、对应题号。\par
          \vspace{0.8mm}
          第八章另附套卷未考、但与这些工具直接相连的结论与定理。\par}
      \end{minipage}
    };
    \node[
      anchor=south west,
      text=Muted,
      font=\sffamily\small,
      align=left
    ] at ([xshift=24mm,yshift=22mm]current page.south west) {
      供同类测评复习 \quad\(\cdot\)\quad 侧重小学课本之外的一层
    };
  \end{tikzpicture}%
  \null\clearpage
}

\newcommand{\fact}[1]{%
  \begin{tcolorbox}[
    enhanced,
    colback=Mint,
    colframe=Teal!38,
    boxrule=0.45pt,
    arc=2pt,
    left=3.2mm,right=3.2mm,top=1.6mm,bottom=1.6mm,
    before skip=1.6mm,
    after skip=1.2mm
  ]
    #1
  \end{tcolorbox}%
}

\newcommand{\seen}[1]{%
  \par\vspace{0.15em}
  {\sffamily\small\color{TealDark}卷面}\hspace{0.55em}{\small #1}\par
}

\newcommand{\linkto}[1]{%
  \par\vspace{0.15em}
  {\sffamily\small\color{TealDark}承接}\hspace{0.55em}{\small #1}\par
}

\newcommand{\ext}[1]{%
  \par\vspace{0.05em}
  {\sffamily\small\color{Coral!85!black}拓展}\hspace{0.55em}{\small #1}\par
}

\newcommand{\chap}[2]{%
  \Needspace{7.2cm}%
  \phantomsection
  \addcontentsline{toc}{section}{\texorpdfstring{#1\hspace{0.6em}#2}{#1 #2}}%
  \sectionbanner{#1}{#2}%
}

\newcommand{\topic}[1]{%
  \Needspace{4.6cm}%
  \phantomsection
  \addcontentsline{toc}{subsection}{#1}%
  \par\vspace{1.6mm}%
  {\sffamily\bfseries\color{TealDark}#1}\par
  \vspace{0.5mm}%
  {\color{Teal!45}\hrule height 0.45pt}%
  \vspace{1.3mm}%
}

\makeatletter
\renewcommand*\l@section[2]{%
  \ifnum \c@tocdepth >\z@
    \addpenalty\@secpenalty
    \addvspace{0.38em \@plus 0.4pt}%
    \setlength\@tempdima{1.4em}%
    \begingroup
      \parindent \z@ \rightskip \@pnumwidth
      \parfillskip -\@pnumwidth
      \leavevmode \bfseries
      \advance\leftskip\@tempdima
      \hskip -\leftskip
      #1\nobreak\hfil
      \nobreak\hb@xt@\@pnumwidth{\hss #2}\par
    \endgroup
  \fi}
\renewcommand\tableofcontents{%
  \sectionbanner{目录}{先按类别，再进入具体工具}%
  \begingroup
  \small
  \linespread{0.92}\selectfont
  \@starttoc{toc}%
  \par\vspace{2.2mm}
  {\sffamily\color{Muted}一题若同时用到两类工具，表中记在更关键的一类；正文相应条目仍会点到另一类。\par}%
  \endgroup
}
\makeatother

\begin{document}
\raggedbottom
\outlinecover

\tableofcontents
\clearpage

\chap{导读}{三套卷在考什么}

这三套卷的题面是行程、时钟、骰子、灯和图形，底下反复用的是同一层工具：先分类或找一个不变的量，再计数或递推，最后用余数、因数或对称性收尾。小学课本里的四则运算、分数、百分数、平均数、长方形和圆，在卷里只作为语言，下面不再重复。

前七章写课本停住之后、三套卷真正拿来用的那一层。每条先给定义，再给一条可以记住的事实，并注明它出现在哪一套的哪些题。标有「拓展」的，是同一条事实的紧邻一步。第八章单独列出套卷没有考到的定理和公式，题号表不收录。

\begin{tcolorbox}[
  enhanced,
  colback=Sky,
  colframe=Teal!35,
  boxrule=0.5pt,
  arc=2.5pt,
  left=3.5mm,right=3.5mm,top=2.4mm,bottom=2.4mm
]
{\sffamily\bfseries\color{TealDark}记号}\par\vspace{1mm}
\(n!\) 表示 \(1\times2\times\cdots\times n\)。
\(P_n^k\) 表示从 \(n\) 个不同元素里取 \(k\) 个的排列数，
\(C_n^k\) 表示相应的组合数。
\(\varphi(n)\) 表示 \(1\) 到 \(n\) 中与 \(n\) 互质的正整数个数。
\end{tcolorbox}

\vspace{1mm}
{\sffamily\bfseries\color{TealDark}题号落点}\par
\vspace{1mm}

{\small\setlength{\tabcolsep}{3.2pt}%
\noindent
\begin{tabularx}{\linewidth}{@{}>{\raggedright\arraybackslash\sffamily\bfseries}p{28mm}>{\raggedright\arraybackslash}X>{\raggedright\arraybackslash}X>{\raggedright\arraybackslash}X@{}}
\toprule
类别 & 第三套 & 第四套 & 第五套 \\
\midrule
计数 & 6、10、15、16、21、22 & 1、17 & 2、9、13 \\
概率 & 5 & 16 & 8、11 \\
递推与周期 & 1、7 & 20 & 5、12、18 \\
整除与因数 & 2、3、12 & 5、11、15、21 & 7、14、16 \\
等式与最值 & 9、11、14、19 & 2、3、6、8、10、12、14 & 1、4、10 \\
图形与运动 & 4、8、18 & 4、7、9、13 & 3、6、15、17 \\
操作与不变量 & 13、17、20 & 18、19、22 & --- \\
\bottomrule
\end{tabularx}
}

\vspace{2mm}
多选题里的「一定正确」，要对所有允许的情况成立；「可以」只要有一个例子。判断之前，先把题目里的不变量或全部可能写出来，再逐句对照。

\chap{一、计数}{分类、分步、有没有顺序}

「有多少种」先分清三件事：是分成几类还是分成几步，结果算不算顺序，以及有没有成对出现、可以只算一半的情形。

\topic{加法原理与乘法原理}

一件事有几类做法，类与类互不重叠，总数是各类方法数之和。一件事要按顺序做完几步，总数是各步方法数之积。

\fact{分类用加法，分步用乘法。先写清「类」或「步」的界限，再计算。}

两位偶数：个位先在偶数里选，十位再在剩下的卡片里选，这是分步。棋子可以全放黑格，也可以全放白格，两类相加。山路上的走法分成「到达必经点」和「离开必经点」两段，两段的方法数相乘。

\seen{第三套第 21、22 题；第五套第 2、13 题。}

\topic{排列与组合}

排列计顺序，组合不计顺序。从 \(n\) 个不同元素中取 \(k\) 个：

\fact{\(P_n^k=n(n-1)\cdots(n-k+1)\)，\quad \(C_n^k=\dfrac{P_n^k}{k!}\)。}

数字排成一列用排列。从 \(1\) 到 \(12\) 里取两个数、不计顺序，用组合。拿不准时先问：交换两个位置后，算不算新的一种。

\seen{第三套第 10、15 题；第五套第 2 题。}

\topic{对称、补集与容斥}

全体结果里，若「严格大于」和「严格小于」一一对应，就只保留一半，并单独减去双方持平的情形。

\fact{对称时，所求 \(=\dfrac{\text{全部}-\text{持平}}{2}\)。\\
「至少一个」\(=\text{全部}-\text{一个都没有}\)。\\
两类都算上时，\(|A\cup B|=|A|+|B|-|A\cap B|\)。}

四个不同数字排成一列，比较前两位之和与后两位之和，大于和小于对称，等于要另数。格子里「没有公共点的两格」，用全部配对减去共边或碰角的配对。

\seen{第三套第 10、16 题；第五套第 11 题。}
\ext{两根指针各自落在某段区间，问至少一根落在其中，就是上面的并集公式。第四套第 9 题的喷泉用到它。}

\topic{按余数分类}

两个整数的和能被 \(m\) 整除，等价于它们除以 \(m\) 的余数之和等于 \(0\) 或 \(m\)。先把数按余数分组，再在组内、组间分别用组合数，最后把各类相加。

\fact{两个数的和能被 \(3\) 整除：两个余数都是 \(0\)，或一个是 \(1\)、一个是 \(2\)。都是 \(1\) 时余数之和为 \(2\)，都是 \(2\) 时余数之和为 \(4\)，这两种都不能被 \(3\) 整除。}

\seen{第三套第 15 题。}
\ext{乘积能被某个数整除时，同样先按余数分组，只是条件改成余数之积满足要求。概率题里的骰子经常这样预处理。}

\topic{抽屉原理}

要保证出现某种情况，先构造一种还没出现该情况的取法，并且让这种取法尽量多，再多取 \(1\) 个。

\fact{保证「至少一种达到 \(a\) 个，并且另有一种达到 \(b\) 个」时，最不利的是：一种颜色可以很多，其余颜色都停在 \(b-1\) 个。}

\seen{第四套第 1 题。}
\ext{\(n\) 个物体放进 \(k\) 个抽屉，必有一个抽屉至少有 \(\lceil n/k\rceil\) 个。}

\topic{图形与结构的分类}

数三角形、数路径、数两人小组，都是同一句话：分类标准要互斥，并且合起来盖住全部。

两人小组若要求任意两组都有一名公共学生，结构只有两类：所有小组共用同一名学生；或者恰好三个小组围成一个三角形。计数时按这两类分别做，再相加。

\fact{先列出互斥的类型，再在每一类里用乘法原理。}

\seen{第三套第 6、22 题；第四套第 17 题。}

\chap{二、概率}{先定样本空间，再缩小范围}

三套卷里的概率都是古典概型：每个结果可能性相同。条件概率是在此基础上把分母换成「条件已经发生」的那些结果。

\topic{古典概型}

\fact{\(P(A)=\dfrac{|A|}{|\Omega|}\)。\(\Omega\) 是全部等可能结果，\(|A|\) 是其中属于事件 \(A\) 的个数。}

动手前先写明一件事：怎样算「一种结果」。两枚骰子有 \(6\times6=36\) 种。球除颜色外可区分时，按组合数计数；题目若声明除颜色外完全相同，则按颜色组成计数。两种约定不能混在同一次计算里。

从若干个数里任取两个，问乘积满足某条件的概率，分子分母都用组合数。

\seen{第五套第 8、11 题；第三套第 21 题是同一套计数，只问方法数。}

\topic{条件概率}

已知事件 \(B\) 已经发生，再求 \(A\) 的概率，分母改成 \(B\) 中的结果数。

\fact{\(P(A\mid B)=\dfrac{|A\cap B|}{|B|}\)。}

两枚骰子已知点数之积能被 \(4\) 整除，再求点数之和为偶数的概率：先数出积能被 \(4\) 整除的掷法，只在这些掷法里数和为偶数的。分母就是这些掷法的个数。

不放回地先取两个、再取第三个，并且已知第一次至少有一个黄球，也是同一条：先固定「第一次至少一黄」的那些完整取法，再在其中数所问的事件。每条完整取法要按同一种规则书写，这样它们才等可能。

\seen{第三套第 5 题；第四套第 16 题。}
\ext{几步彼此独立时（有放回，或两枚骰子），概率可以相乘。不放回时前后相关，把整条取法当作一个结果来数。}

\topic{先分组，再把组的大小平方}

四枚骰子分成红绿、蓝黄两组，问两组乘积相等的掷法有多少。设某个乘积 \(P\) 恰好有 \(n\) 种掷法，则两组同时得到 \(P\) 的方法有 \(n^2\) 种。对所有可能的 \(P\) 把 \(n^2\) 加起来。

\fact{按结果的值分组。组的大小为 \(n\) 时，两组同时落入该组有 \(n^2\) 种。}

这仍是计数。它说明：看起来像概率的综合题，常常先做成「每种结果有几种」。

\seen{第三套第 21 题。}

\chap{三、递推与周期}{状态会重复，就不要逐项硬算}

递推题问的通常是很后面的某一项、某一项的余数，或者到达某一格的方法数。状态种类有限时，序列一定进入循环。

\topic{循环节}

只依赖前一项或前两项的规则，若只关心奇偶或除以 \(m\) 的余数，可能的状态只有有限个。状态一重复，后面就循环。

\fact{把奇偶或余数写出来，直到一段状态与开头重复。循环节的长度就是这两处之间的步数。}

后一项等于前两项之和时，奇偶往往每三项一组。后一项等于前一项的 \(2\) 倍再加 \(1\) 时，除以 \(7\) 的余数最多 \(7\) 种，循环节不超过 \(7\)。前 \(30\) 项里的偶数个数、第 \(20\) 项与第 \(21\) 项的余数，都落在循环节的确定位置上。

\seen{第三套第 1、7 题。}
\ext{只依赖前一项余数的递推，模 \(m\) 的循环节长度不超过 \(m\)。依赖前两项时，状态是一对余数，长度不超过 \(m^2\)。}

\topic{步数递推与必经点}

每次向前走 \(1\) 步、\(2\) 步或 \(3\) 步，设 \(f(n)\) 为恰好走到第 \(n\) 格的方法数。走到 \(n\) 的最后一步，来自 \(n-1\)、\(n-2\) 或 \(n-3\)。

\fact{\(f(n)=f(n-1)+f(n-2)+f(n-3)\)，并令 \(f(0)=1\)，下标为负时 \(f=0\)。\\
中途必须停在第 \(k\) 格时，方法数为 \(f(k)\times f(n-k)\)。}

起点到终点被必经点切开，两段独立，所以相乘。卡片放回后再抽，每一步的可选步长不变，上面的递推可以直接用。

\seen{第五套第 18 题。}

\topic{格子路径}

只能向右或向下，且障碍不能经过时，每个格子的方法数来自左边一格和上边一格。

\fact{当前格的方法数 \(=\) 左格的方法数 \(+\) 上格的方法数。障碍格记 \(0\)。按从左上到右下的顺序填表。}

这和上一条是同一件事：当前位置的方法数，等于能一步到达它的那些位置的方法数之和。

\seen{第五套第 9 题。}
\ext{若允许再多一个方向，仍然是「前驱之和」，只是每个格子的前驱变多。有回路时这个填表顺序失效，三套卷里的路径都没有回路。}

\topic{周期位置与螺旋}

一组图形按固定顺序反复出现时，第 \(n\) 个图形对应周期里的第
\[
((n-1)\bmod T)+1
\]
个，其中 \(T\) 是一组的个数。

顺时针螺旋填正整数时，先定圈层，再沿圈数。

\fact{中心单独是第 \(0\) 圈。第 \(k\) 圈（\(k\ge 1\)）有 \(8k\) 个数，这一圈的最大数是 \((2k+1)^2\)。\\
数 \(N>1\) 所在的圈，是使 \((2k+1)^2\ge N\) 的最小正整数 \(k\)。}

圈层确定以后，从该圈的起点沿题目规定的方向逐边数到 \(N\)，就能读出它的邻居和同排、同列。

\seen{第五套第 5 题；第四套第 20 题。}

\topic{长和式的余数}

\(n\) 足够大以后，\(n!\) 会包含模数的全部质因数，从而 \(n!\equiv 0\)。后面再加多少项，余数都不变。

\fact{除以 \(15\) 时，从 \(5!\) 起每一项都能被 \(15\) 整除。\\
\(1!+2!+\cdots+2026!\) 的余数等于 \(1!+2!+3!+4!\) 的余数。}

\seen{第五套第 12 题。}
\ext{模数换成 \(100\) 时，看从哪一项起同时有足够多的因子 \(2\) 和 \(5\)。问法变了，截断的位置跟着变，做法不变。}

\chap{四、整除与因数}{先分解质因数，再数指数}

课本会做短除和最大公因数。这三套卷往前多走的是：指数可以自由选取的范围、互质的个数，以及一个乘积里某个质因数出现几次。

\topic{最大公因数与最小公倍数}

若干盏灯再次同时亮起的间隔，是它们周期的最小公倍数。

若 \(\gcd(a,b)=d\)，则 \(a=da_1\)，\(b=db_1\)，且 \(\gcd(a_1,b_1)=1\)。和固定时，先做这个替换，再在 \(a_1\) 与 \(b_1\) 互质的条件下求积的最大。

\fact{\(\gcd(a,b)\times\operatorname{lcm}(a,b)=ab\)。\\
多个数时，最小公倍数的每个质因数指数，取这些数里的最大指数。}

圆圈上每个数是相邻两数的最大公因数或最小公倍数时，把每个质因数的指数单独看：取最大公因数就是取较小指数，取最小公倍数就是取较大指数。指数沿圆圈要能形成一座先升后降或先降后升的山。

\seen{第三套第 2、3 题；第四套第 15 题。}

\topic{因数的个数}

设 \(N=p_1^{a_1}p_2^{a_2}\cdots p_k^{a_k}\)，质因数两两不同。

\fact{正因数个数为 \((a_1+1)(a_2+1)\cdots(a_k+1)\)。\\
正因数个数为奇数，当且仅当 \(N\) 是完全平方数。}

由此直接得到几条常用的变形：

\begin{itemize}
  \item 奇数因数：去掉因子 \(2\)，对其余指数再用上面的公式。
  \item 能被 \(15=3\times5\) 整除的因数：\(3\) 和 \(5\) 的指数至少为 \(1\)，其余指数仍可自由选取。
  \item 完全平方因数：每个指数只能取 \(0,2,4,\ldots\) 中不超过原指数的那些值。
\end{itemize}

\(1\) 到 \(200\) 里因数个数为偶数的数，个数等于 \(200\) 减去其中完全平方数的个数。

\seen{第三套第 12 题；第五套第 7 题。}

\topic{末尾的 0 与整除次数}

乘积末尾连续 \(0\) 的个数，等于乘积中因子 \(10\) 的个数，也就是质因数 \(2\) 与 \(5\) 能配成的对数。连续整数相乘时，\(2\) 通常更多，于是由 \(5\) 的个数决定。

问一个整数能被 \(5\) 整除多少次，就是数它的质因数分解里 \(5\) 的指数。\(25\) 贡献两个，\(125\) 贡献三个。

\fact{末尾 \(0\) 的个数 \(=\min(\text{因子 \(2\) 的个数},\text{因子 \(5\) 的个数})\)。\\
整除次数 \(=\) 该质因数的指数。}

若干个分数的分子全部相乘时，先把每个分子里的 \(5\) 数出来再相加；若式子里还有分母，分母中的 \(5\) 要减掉。

\seen{第四套第 5、21 题。}

\topic{互质与欧拉函数}

\(\varphi(n)\) 是 \(1,2,\ldots,n\) 中与 \(n\) 互质的正整数个数。不超过 \(n\) 且与 \(n\) 互质，问的就是 \(\varphi(n)\)。

\fact{若 \(n=p_1^{a_1}\cdots p_k^{a_k}\)，则
\(\varphi(n)=n\Bigl(1-\dfrac{1}{p_1}\Bigr)\cdots\Bigl(1-\dfrac{1}{p_k}\Bigr)\)。}

先把 \(n\) 分解质因数，再代入。公式只依赖 \(n\) 有哪些质因数，同一质因数的更高次方已经由前面的 \(n\) 承担。

最简真分数按分母分组时，分母为 \(m\)（\(m\ge 2\)）的一组，分子是 \(1\) 到 \(m-1\) 里与 \(m\) 互质的数，因此恰好有 \(\varphi(m)\) 个。三套卷把分母为 \(2\) 的称作第 \(1\) 组，所以第 \(n\) 组的分母是 \(n+1\)，个数是 \(\varphi(n+1)\)。

分子固定为 \(3\) 时，还要加上 \(3\) 小于分母，以及 \(3\) 与分母互质。

\seen{第五套第 14 题；第四套第 21 题。}

\topic{分数的夹逼}

设 \(\dfrac{p}{q}<\dfrac{r}{s}\)，并且交叉差 \(|ps-qr|=1\)。这时这两个分数挨得最近：夹在它们中间的任何分数，分母都至少是 \(q+s\)。

\fact{交叉差的绝对值为 \(1\) 时，分母最小的中间分数是 \(\dfrac{p+r}{q+s}\)。}

\(\dfrac{3}{8}\) 与 \(\dfrac{2}{5}\) 的交叉差绝对值是 \(1\)，所以满足不等式的分数，分母最小可以是 \(13\)。

\seen{第四套第 11 题。}

\topic{单位分数的拆分}

分子为 \(1\) 的分数可以裂成两个分母不同的单位分数：

\fact{\(\dfrac{1}{n}=\dfrac{1}{n+1}+\dfrac{1}{n(n+1)}\)。}

古埃及分数要求每一项分子都是 \(1\)，并且分母互不相同。分子大于 \(1\) 时，可以先写成若干个相同单位分数之和，再对重复的分母用上面的等式裂开，直到分母全都不相同。

\seen{第五套第 16 题。}
\ext{同一个分数的埃及分数写法往往不止一种。题目若规定项数或分母从小到大，先按项数最少去试，再用这条等式调整重复的分母。}

\chap{五、等式与最值}{先消元，再看极端}

这一类在课本里已有方程和平均数。卷面多出来的是：几个对称式互相换算、和固定时积最大、以及带绝对值的和在何处最小。

\topic{和、积与比}

\fact{\((a+b)^2=a^2+2ab+b^2\)，\quad \(\dfrac{a}{b}+\dfrac{b}{a}=\dfrac{a^2+b^2}{ab}\)。}

已知 \(a+b\) 和 \(\dfrac{a}{b}+\dfrac{b}{a}\)，先求出 \(a^2+b^2\)，再求出 \(ab\)。两个正数的和固定时，它们越接近，积越大；都是正整数时，最大积出现在两数相等或相邻。

三个连续正整数可设为 \(n-1,n,n+1\)。它们的和是 \(3n\)，积是 \(n(n^2-1)\)，公有的 \(n\) 可以约去。

\(a+b\) 为偶数时，\(a\) 与 \(b\) 同奇偶，\(a^2+b^2\) 为偶数。

\seen{第三套第 9、14 题；第四套第 6 题。}

\topic{不定方程与不等式}

二元一次不定方程先让系数变小：各项有公因数就先约去，再枚举较小的那个未知数。两种物品都买、单价和总价都是整数时，枚举的范围是有限的，把正整数解列全，再判断哪一句话对所有解成立。

两个同向的不等式可以相加。由
\[
  a+b>c+d,\qquad b+d>a+c
\]
相加得到 \(2b>2c\)，即 \(b>c\)。

\fact{不定方程：约简后枚举。\\
不等式：同向才能相加；相加后仍然成立的关系，才是「一定」成立的。}

\seen{第三套第 11 题；第四套第 2、3 题。}

\topic{平均、速度与年龄}

总分等于人数乘平均分。把若干人从一组调到另一组时，列的是「被调走的那些人的分数之和」，两组的人数同时改变。

两段路程相等、速度不同时，全程平均速度是

\fact{\(v=\dfrac{2v_1v_2}{v_1+v_2}\)。}

年龄题设今年的年龄，把「若干年前」「若干年后」写成两个等式。行程题先画方向：谁在走、谁已返回、相遇时两边路程之间的关系，再列方程。百分数先还原成「部分 \(=\) 总数 \(\times\) 百分数」，转入或转出之后总数要重算。

\seen{第四套第 2、8、12 题；第五套第 1、4 题。}

\topic{乘积最大与绝对值之和}

把正整数写成若干个正整数之和，使这些加数的乘积最大：加数只取 \(2\) 和 \(3\)，并且 \(3\) 尽量多。若按连排 \(3\) 之后剩下 \(1\)，就把一个 \(3\) 和这个 \(1\) 改成两个 \(2\)，因为 \(3\times1<2\times2\)。

\fact{最优拆分只含 \(2\) 和 \(3\)；剩下 \(1\) 时，把 \(3+1\) 改成 \(2+2\)。}

另一类最小值是带权重的绝对值之和
\[
  F(n)=\sum_{k=1}^{m} k\,|n-k|.
\]
位置 \(k\) 上的权重是系数 \(k\)。从左向右把权重累加，累加和第一次达到总权重一半的那个位置，就是最小值点。权重全为 \(1\) 时，这个位置就是中位数。

\seen{第五套第 10 题；第四套第 10 题。}

\topic{在约束下取极端}

和、个数、最小值都已固定，又要求互不相同的正整数，要使最大数尽量小：其余的数取满足条件的最小可用整数，余下的部分交给最大数。

五个数的十个两两正差互不相同：差的个数是 \(C_5^2=10\)，这 \(10\) 个差至少要占 \(10\) 个不同的正整数。先用这个估计下界，再构造一组尽量紧的数，逐一检查差是否重复。

\fact{「尽量小」的最大值：其余量先取最小可用值。\\
「互不相同的差」：先数差有几个，再构造并检查。}

\seen{第三套第 19 题；第四套第 14 题。}

\chap{六、图形与运动}{比例、分段、展开}

图形题在这三套卷里很少要求添加辅助线大全套。常见的是：把长度或面积化成比，把运动分成几段，以及把反射改成直线。

\topic{长度与面积的比}

若干个圆排成一排时，总长度等于各直径之和，再减去相邻圆重叠掉的部分。题目若给出两个端点之间的水平距离，先把「露出」和「重叠」都换成端点之间的距离。

凹多边形的周长，常常等于把它凹进去的边平移出去之后、外接矩形的周长。成对的内部线段方向相反，平移后并到外圈。

正方形被直线分成几块、已知三块求第四块时，先找共高或可以拼成同一个三角形的两块。同高时，面积比等于底之比。

\fact{同高（或同底）的两块，面积比等于底（或高）之比。\\
总长度 \(=\) 各段长度之和 \(-\) 重叠。}

圆之间的梯形，先把上底、下底、高写成半径的倍数，再代梯形面积公式。

\seen{第三套第 4、18 题；第四套第 7、13 题；第五套第 6 题。}

\topic{动点图像与时钟}

动点匀速沿折线移动时，三角形面积随时间的图像是分段直线。某一段里底和高都不变，图像水平；其中一量均匀变化，图像是斜线段。先按拐弯把路线分段，再写每段起点和终点的面积。

时钟上，分针每分钟走 \(6^\circ\)，时针每分钟走 \(0.5^\circ\)。

\fact{分针与时针的相对速度是每分钟 \(5.5^\circ\)。\\
夹角为 \(\theta\) 时，\(|30H-5.5m|\) 等于 \(\theta\) 或 \(360^\circ-\theta\)。}

问一根指针有多少时间落在某两个数字之间，先把一小时分成与钟面比例相同的时间，再乘以小时数。两根指针都要计入时，用并集，重叠的时段只算一次。

\seen{第三套第 8 题；第四套第 4、9 题。}

\topic{涂色与卷筒}

堆叠的单位正方体涂外表面，并且与地面接触的底面不涂。一个小块被涂到，是因为它有一面露在空气中。完全未涂色的块，是每一面要么贴着别的块、要么只贴着地面。

\fact{按层看：侧面最外一圈、顶上露出来的一面会被涂到。\\
只贴地面、其余各面都被挡住的块，算未涂色。}

卷筒卫生纸的纸料体积，等于整卷圆柱的体积减去纸筒的体积。用这个体积除以纸的厚度和纸宽，得到纸长，再除以每节的长度。包装上的层数要乘进单层厚度。长度、半径的单位先化成同一套。

\seen{第五套第 3、15 题。}

\topic{反射等于展开}

点在多边形内直线前进，碰到边时入射角等于反射角，碰到顶点就停止。把多边形沿该边镜像出去之后，反射后的路径变成镜像里的一条直线。

\fact{反弹一次，相当于进入相邻的一张镜像。\\
直线穿过第几条像边，就是第几次碰到的边；直线穿过像顶点，运动在对应顶点停止。}

正六边形可以铺满平面，所以连续几次反弹对应这条直线依次穿过的那些像。第一次碰到指定的边，用来限定直线的方向；再要求若干次反弹后到达某个顶点，就在展开图里量对应的线段。

\seen{第五套第 17 题。}

\chap{七、操作与不变量}{每次操作都留下一个改不掉的量}

操作题若逐步模拟，步骤会很多。先找一个操作改变不了的量，用来说明哪些局面到不了；若每一步都影响后面、而不影响前面，就从小到大唯一确定。

\topic{不变量}

\fact{操作前后保持不变的量，用来排除到不了的局面。\\
会变、但变化方式固定的量，用来缩小可能的局面。}

三堆棋子每次「一堆加 \(2\)、另两堆各减 \(1\)」：总数不变。被减的两堆奇偶同时翻转，奇数堆的个数因此保持偶数或保持奇数。

黑板上每次擦去两个数、写上它们的和：最后剩下的数等于原来所有数的总和。

数对 \((x,y)\) 变成 \((x+y,y)\) 或 \((x,x+y)\) 时，
\[
  \gcd(x+y,y)=\gcd(x,y).
\]
从 \((1,1)\) 出发，以后每一对的最大公因数都是 \(1\)。若两个坐标不相等，前驱只有一个：从较大的数里减去较小的数。能否恰好 \(k\) 步到达，就看这样倒推 \(k\) 次是否回到 \((1,1)\)。

\seen{第三套第 13、17 题；第四套第 18 题。}

\topic{开关与从小到大}

按下一盏灯，它自己改变亮灭；与它相邻的灯里，只有数字更大的那些一起改变。数字更小的灯不受这次操作影响。同一盏灯按两次，等于没有按。

\fact{影响只指向编号更大的灯。\\
从编号小的灯开始：它现在的状态与目标不同，就必须按它；决定后不再回头。}

最终亮灭只与每盏灯被按的次数是奇数还是偶数有关。题目要的按键序列，可以写成这些灯的编号从小到大。

\seen{第三套第 20 题。}
\ext{若每盏灯影响的是一个固定的集合，并且可以排出「先决定的灯不再被后面的操作改变」的顺序，就都能这样消元。排不出这种顺序时，再把每次按下看成模 \(2\) 的一次加法。}

\topic{天平的三种结果}

无砝码天平一次称量有三种结果：左重、右重、平衡。只知道其中恰有一枚偏重时，一次称量把可能的硬币分成三类。

\fact{称 \(k\) 次，最多区分 \(3^k\) 枚里面的哪一枚偏重。\\
\(2\) 次最多区分 \(9\) 枚。}

第一次把 \(3\) 枚对 \(3\) 枚。平衡，则重的在未上秤的 \(3\) 枚里；不平衡，则重的在较重那一盘的 \(3\) 枚里。第二次在剩下的 \(3\) 枚里取 \(1\) 枚对 \(1\) 枚。\(10\) 枚超出 \(3^2\)，不能保证两次找出。

\seen{第四套第 19 题。}
\ext{若不知道偏重还是偏轻，每种硬币对应两种可能，\(k\) 次最多处理 \((3^k-1)/2\) 枚。题目先写明「已知偏重」还是「轻重未知」。}

\topic{圆盘与弧长}

几个圆盘相切、转动时不打滑，接触点走过的弧长相等。相邻圆盘的转向相反。

\fact{弧长相等，所以转过的角度与直径成反比。\\
直径为 \(d_1,d_2\) 时，转角满足 \(\theta_1 d_1=\theta_2 d_2\)。}

刻度把圆周分成 \(24\) 等份，每份 \(15^\circ\)。先由直径比算出转角，再换成走过几格，最后按该盘的转向确定正上方的新数字。

侧面涂色被印到下一只盘，印上的也是相等的弧长。要让一只盘的侧面全部被印到，传到它上面的弧长总量至少等于它的周长；颜色只涂在一段弧上时，数需要经过几个这样的涂色段。

\seen{第四套第 22 题。}

\chap{八、邻近定理}{套卷未考，与前七章直接相连}

下面每条都不对应三套卷里的某一道题。它们是前七章已经在用的工具，再往前得到的结论、公式或定理。用的时候先回到括号里的那一节，确认问法已经从「数一次」变成了这里的形式。

\topic{组合数恒等式}

从 \(n\) 个不同元素里取 \(k\) 个，与留下 \(n-k\) 个是同一件事。杨辉三角里每个数等于肩上两数之和。\(n\) 元集合的全部子集，是对每个元素选「进或不进」。

\fact{\(C_n^k=C_n^{n-k}\)。\\[0.4em]
\(C_n^k=C_{n-1}^{k}+C_{n-1}^{k-1}\)（\(n\ge 1\)，约定越界的组合数为 \(0\)）。\\[0.4em]
\(\displaystyle\sum_{k=0}^{n}C_n^k=2^n\)。\\[0.3em]
\(n\ge 1\) 时，\(k\) 为偶数的各项之和与 \(k\) 为奇数的各项之和都等于 \(2^{n-1}\)。}

二项式定理把同一批系数用在乘方上：
\[
(a+b)^n=\sum_{k=0}^{n}C_n^k\,a^{n-k}b^k.
\]

\linkto{第一章 · 排列与组合。问「取 k 个」与「排除 k 个」、或要给组合数找递推时用。}

\topic{三个集合的容斥}

两个集合的并，前七章已经写成加起来再减去交。三个条件同时出现时，两两的交会被减两次，三个一起的部分要加回去。

\fact{\(|A\cup B\cup C|
=|A|+|B|+|C|
-|A\cap B|-|A\cap C|-|B\cap C|
+|A\cap B\cap C|\)。}

「至少满足一条」先算并集。「一条都不满足」用全部减去这个并集。

\linkto{第一章 · 对称、补集与容斥。两个集合的公式是上式只留 \(A,B\) 的情形。}

\topic{全概率公式}

先把全部结果分成几类，类与类不重叠、合起来盖住全部。事件 \(A\) 的概率等于它在每一类里的条件概率，乘上该类的概率，再相加。

\fact{若 \(B_1,\ldots,B_m\) 两两不交，并且并起来是全部结果，则
\[
P(A)=\sum_{i=1}^{m}P(A\mid B_i)\,P(B_i).
\]
\(A\) 与 \(B\) 独立时，\(P(A\cap B)=P(A)P(B)\)，从而 \(P(A\mid B)=P(A)\)。}

分类标准要互斥。不放回抽样前后相关，不能把两步的概率直接相乘。

\linkto{第二章 · 条件概率。条件概率是已知落在某一类之后再算；全概率是先分类、再把各类加权加回去。}

\topic{同余运算与欧拉定理}

两个数除以 \(m\) 的余数相同，就说它们模 \(m\) 同余。同余可以先合并再取余数，也可以先取余数再合并，结果相同。

\fact{若 \(a\equiv b\pmod m\)，\(c\equiv d\pmod m\)，则
\[
a+c\equiv b+d,\qquad ac\equiv bd,\qquad a^k\equiv b^k\pmod m.
\]
若 \(\gcd(a,n)=1\)，则 \(a^{\varphi(n)}\equiv 1\pmod n\)。\\
因此 \(a^k\bmod n\) 可以由指数模 \(\varphi(n)\) 决定。\\
\(n=p\) 为素数且 \(p\) 不整除 \(a\) 时，\(a^{p-1}\equiv 1\pmod p\)。}

末位、末两位就是模 \(10\)、模 \(100\)。\(\gcd(a,n)>1\) 时欧拉定理不能用，例如 \(2\) 的方幂的末位要单独看循环。

\linkto{第四章 · 互质与欧拉函数；第三章 · 循环节。\(\varphi(n)\) 从「个数」变成「指数可以缩到的模」。}

\topic{阶乘里的质因数指数}

\(n!\) 里质因数 \(p\) 的个数，等于 \(1\) 到 \(n\) 中 \(p\) 的倍数、\(p^2\) 的倍数、\(p^3\) 的倍数……各自贡献的次数之和。

\fact{\(n!\) 中质因数 \(p\) 的指数为
\[
\left\lfloor\frac{n}{p}\right\rfloor+\left\lfloor\frac{n}{p^2}\right\rfloor+\left\lfloor\frac{n}{p^3}\right\rfloor+\cdots
\]
加到 \(p^k>n\) 为止。\\
\(n!\) 末尾连续 \(0\) 的个数，等于其中质因数 \(5\) 的指数。}

连续若干个数的乘积，若长度不长，仍逐个数分解；连乘到 \(n!\) 时用上面的和，避免把每一项都写出来。

\linkto{第四章 · 末尾的 \(0\) 与整除次数。短乘积数 \(2\) 和 \(5\)；阶乘换成这条求和。}

\topic{裴蜀定理}

两个整数的线性组合 \(ax+by\)，系数允许为正、为负或为零时，能得到的全部整数，恰好是 \(\gcd(a,b)\) 的倍数。

\fact{\(ax+by=c\) 有整数解，当且仅当 \(\gcd(a,b)\) 整除 \(c\)。\\
若还要求 \(x,y\) 都是非负整数，并且 \(\gcd(a,b)=1\)，则不能表示的最大正整数是 \(ab-a-b\)。\\
\(\gcd(a,b)>1\) 时，不能被这个最大公因数整除的数都无法表示，因此没有「最大的一个」。}

求一组整数解时，用辗转相减把系数缩小，这和数对 \((x,y)\to(x+y,y)\) 是互逆的两步。

\linkto{第五章 · 不定方程；第七章 · 不变量。套卷里的方程系数已约简、且解都为正，枚举就够；这里给出有没有解、以及非负解的边界。}

\topic{中国剩余定理}

同时满足两个余数条件，先看两个余数是否相容，再看解隔多少个数重复一次。

\fact{方程组
\[
x\equiv a\pmod m,\qquad x\equiv b\pmod n
\]
有解，当且仅当 \(a\equiv b\pmod{\gcd(m,n)}\)。\\
有解时，解在模 \(\operatorname{lcm}(m,n)\) 下唯一。\\
若 \(\gcd(m,n)=1\)，则一定有解，并且模 \(mn\) 唯一。此时若 \(mu\equiv 1\pmod n\)，则
\[
x\equiv a+mu(b-a)\pmod{mn}.
\]}

多个两两互质的模，解在这些模的乘积下唯一。\(u\) 用辗转相减来找：它就是 \(m\) 模 \(n\) 的乘法逆元。

\linkto{第四章 · 最小公倍数。三盏灯同时亮起，是每个余数都为 \(0\) 的特殊情形，条件自动相容。余数不全为 \(0\) 时，先用上面的相容条件。}

\topic{等比求和与一阶递推}

公比相同的一串数，前 \(n\) 项可以一次加总。后一项是前一项的 \(p\) 倍再加一个常数时，先找一个不动点，偏差就变成等比数列。

\fact{首项 \(a\)、公比 \(q\) 的前 \(n\) 项和：\(q=1\) 时为 \(na\)；\(q\neq 1\) 时为 \(\dfrac{a(1-q^n)}{1-q}\)。\\[0.5em]
\(a_{n+1}=pa_n+r\)。若 \(p=1\)，则 \(a_n=a_1+(n-1)r\)。\\
若 \(p\neq 1\)，令不动点 \(c=\dfrac{r}{1-p}\)，则
\[
a_n-c=p^{n-1}(a_1-c).
\]}

只问余数时，仍用循环节，不必先写通项。题目要的是第 \(n\) 项本身，或要比较大小时，再用通项。

\linkto{第三章 · 循环节。\(a_{n+1}=2a_n+1\) 在套卷里只取模 \(7\)；通项是同一递推的另一问法。}

\topic{均值不等式}

正数的算术平均不小于几何平均。和固定时，积在各数相等时最大；积固定时，和在各数相等时最小。

\fact{正数 \(a,b\) 满足 \(\dfrac{a+b}{2}\ge\sqrt{ab}\)，等号当且仅当 \(a=b\)。\\[0.45em]
\(n\) 个正数同样：
\[
\frac{a_1+\cdots+a_n}{n}\ge\sqrt[n]{a_1\cdots a_n},
\]
等号当且仅当 \(a_1=\cdots=a_n\)。}

正整数不能随意取成相等的实数。和固定、加数是正整数、要使乘积最大时，仍按第五章把加数取成 \(2\) 和 \(3\)。

\linkto{第五章 · 和、积与比，以及乘积最大。两个正数「越接近积越大」就是这里 \(n=2\) 的情形。}

\topic{染色不变量}

棋盘按黑白相间染色后，每放一块 \(1\times 2\) 的矩形，就同时盖住一黑一白。因此能被这种矩形铺满的区域，黑格数必须等于白格数。

\fact{黑格数 \(\neq\) 白格数时，不能用 \(1\times 2\) 的矩形铺满。\\
从完整棋盘挖去两格后，若这两格同色，剩下的区域也不能铺满。}

这只说明「不能」。黑白格数相等，还要再看区域是否连成可以一块块盖住的形状。

\linkto{第七章 · 不变量；第一章的同色放棋子用的是同一张黑白棋盘，问法从「几种放法」换成「能不能盖住」。}

\topic{相似比}

两个图形对应的长度之比都是 \(k\) 时，面积之比是 \(k^2\)，体积之比是 \(k^3\)。三角形里作一边的平行线，截出的小三角形与原三角形相似。

\fact{长度比 \(k\)，则面积比 \(k^2\)，体积比 \(k^3\)。\\
平行于三角形一边的直线，把另外两边截成相同的比。\\
截出的小三角形与大三角形的面积比，等于对应边之比的平方。}

同高不同底时，面积比仍等于底之比，那是 \(k\) 的一次方，因为高没有同时放大。两条都用之前，先看高有没有变。

\linkto{第六章 · 长度与面积的比。套卷里的面积比多用同高；图形整个放大时改用平方。}

\vspace{6mm}
\begin{tcolorbox}[
  enhanced,
  colback=Sky,
  colframe=Teal!35,
  boxrule=0.5pt,
  arc=2.5pt,
  left=4mm,right=4mm,top=3mm,bottom=3mm
]
{\sffamily\bfseries\color{TealDark}怎样用这份纲要}\par\vspace{1.4mm}
前七章先记住色块里的那一条事实，再按题号回到三套卷，只回答「这题用的是哪一条」。变式若只多一个条件，先看该条下面的「拓展」。问法若换成高次幂的余数、两个余数同时成立、通项、铺砖或相似图形的面积，再到第八章对公式。
\end{tcolorbox}

\end{document}
